\section{来源代码}

华为官方来源：H1 为 2025-09-18 高管主题演讲，H2 为同日 SuperPoD 架构开放稿，H3 为 2026-03-02 MWC 产品稿，H4 为 2026-07-17 Atlas 950 实物系统稿。竞品、产业链和政策官方来源编号为 C1--C28，均记录日期、网址、证据用途与本地原件哈希。

公司来源：F1 为科大讯飞 2025 年年报，F2 为拓维信息 2025 年年报，F3 为申菱环境 2025 年年报，F4 为航天电器 2025 年交易所年报，F5 为航天电器 2026-04-20 投资者问答；其余 12 家建模公司均已归档交易所完整年报原件并核对 FY2025 收入、归母净利与公司行为。卖方部分归档 38 份原始 PDF，覆盖 14 个重点代码和 2 份行业报告；结构化一致预期仅采用正权重原始 PDF。

\section{高影响声明审计摘要}

\begin{longtable}{L{4.0cm}L{2.4cm}L{3.4cm}L{3.4cm}}
\toprule
\textbf{声明} & \textbf{状态} & \textbf{依据} & \textbf{报告采用措辞}\\
\midrule
\endhead
华为第一个超节点 & 拒绝 & 前代 Atlas 900 A3/CloudMatrix384 & Atlas 950 首次公开实物展示\\
1,024 卡实物系统 & 采用 & H4 华为官方 & 物理展示，性能为厂商自述\\
8,192 卡已交付 & 拒绝 & H1/H3 为 Q4 路线图 & 最大设计，等待实物/客户验收\\
256TB HBM & 拒绝 & H4 仅称全局地址空间 & 不标为 HBM\\
申菱为 Atlas 供应商 & 未证实 & F3 仅确认华为客户/H 公司合作 & 一般冷却客户关系、平台订单未确认\\
科大讯飞采购 Atlas & 未证实 & F1 仅确认 950 模型协同 & 下游工作负载需求锚\\
拓维有 Atlas 订单 & 未证实 & F2 为昇腾伙伴和华为购销 & 昇腾生态观察\\
航天电器供货 950 & 未证实 & F4/F5 确认产品能力但没有确认提问前提 & 连接器能力观察\\
四川长虹 AI 服务器 ODM & 排除 & 公司问答明确否认 & 不纳入标的池\\
\bottomrule
\end{longtable}

\section{来源穷尽与剩余缺口}

本案已经检查华为三阶段官方发布、公司年报/摘要/互动问答、原始卖方报告、竞品官方资料和政策文件。仍未找到：Atlas 950 客户与站点、8,192 卡整机验收、系统价格与功耗、完整 BOM、光模块/交换/OCS 数量、CDU/冷板/连接器数量、HBM 制造与封装/良率、OEM/ODM 和上市公司供应商分配、公司订单/收入/毛利。以上项目在来源穷尽日志中均保留下一检索路径。

\section{数据文件索引}

完整来源注册表、逐声明审计、失败和负面检索、市场财务核对、公司关系字段、客户链审计与估值复算日志均作为本报告随附电子证据包保存；正文只展示读者决策所需的来源编号、证据等级与审计结论。
